% TEMPLATE-MILC-v3.tex - plantilla maestra para todas las guias MILC v3
% Ver scripts/generadores/build-guia-milc.py para la lista completa de
% claves esperadas. El script valida que no queden placeholders sin
% reemplazar antes de escribir el archivo final.

\documentclass[11pt,letterpaper]{article}
\usepackage[margin=1.75cm]{geometry}
\usepackage[table]{xcolor}
\usepackage{fontspec}
\usepackage[spanish]{babel}
\usepackage{tikz}
\usepackage[most]{tcolorbox}
\usepackage{tabularx}
\usepackage{array}
\usepackage{fancyhdr}
\usepackage{titlesec}
\usepackage{ragged2e}
\usepackage{enumitem}
\usepackage{microtype}

\setmainfont{Helvetica Neue}
\setsansfont{Helvetica Neue}
\setlength{\parindent}{0pt}
\setlength{\parskip}{6pt}
\hyphenpenalty=200
\exhyphenpenalty=200
\pretolerance=400
\tolerance=2000
\setlength{\emergencystretch}{1em}
\renewcommand{\arraystretch}{1.25}
\setlist{leftmargin=5mm,itemsep=1mm,topsep=1mm}
\newcolumntype{Y}{>{\RaggedRight\arraybackslash}X}

\definecolor{milcMagenta}{HTML}{E6007E}
\definecolor{milcVino}{HTML}{5A0038}
\definecolor{milcTurquesa}{HTML}{0093A5}
\definecolor{milcVerde}{HTML}{58A923}
\definecolor{milcMostaza}{HTML}{E5B400}
\definecolor{milcOcre}{HTML}{B86600}
\definecolor{milcNegro}{HTML}{241816}
\definecolor{milcGris}{HTML}{F7F7F4}
\definecolor{milcLinea}{HTML}{E8E2DD}
\definecolor{milcGradePrimary}{HTML}{84CC16}
\definecolor{milcGradeDark}{HTML}{4D7C0F}
\definecolor{milcGradeSoft}{HTML}{ECFCCB}
\definecolor{milcGradeAccent}{HTML}{CFE7FF}
\definecolor{milcGradeText}{HTML}{FFFFFF}
\color{milcNegro}

\pagestyle{fancy}
\fancyhf{}
\lhead{\small Tecnología e Informática · Grado 7}
\rhead{\small Guía 6 · MILC v3}
\cfoot{\small\thepage}
\renewcommand{\headrulewidth}{0pt}

\newtcolorbox{titlebox}[1]{
  enhanced,colback=#1,colframe=#1,arc=7mm,boxrule=0pt,
  left=7mm,right=7mm,top=3mm,bottom=3mm,
  before skip=8pt,after skip=8pt,
  fontupper=\sffamily\bfseries\Large\color{white}
}

\newtcolorbox{softbox}[3]{
  enhanced,colback=#2,colframe=#1,borderline west={4pt}{0pt}{#1},
  arc=2mm,boxrule=.4pt,left=4mm,right=4mm,top=3mm,bottom=3mm,
  before skip=6pt,after skip=8pt,title={#3},coltitle=white,
  colbacktitle=#1,fonttitle=\sffamily\bfseries,fontupper=\RaggedRight,
  attach boxed title to top left={xshift=3mm,yshift=-2mm},
  boxed title style={arc=2mm, boxrule=0pt}
}

\newtcolorbox{drawbox}[2]{
  enhanced,colback=white,colframe=#1,arc=4mm,boxrule=1pt,
  left=5mm,right=5mm,top=4mm,bottom=4mm,before skip=8pt,after skip=8pt,
  title={#2},coltitle=white,colbacktitle=#1,fonttitle=\sffamily\bfseries,
  fontupper=\RaggedRight,
  attach boxed title to top left={xshift=4mm,yshift=-2mm},
  boxed title style={arc=3mm, boxrule=0pt}
}

\newtcolorbox{infoband}[1]{
  enhanced,colback=#1!8,colframe=#1!50,arc=2mm,boxrule=.5pt,
  left=4mm,right=4mm,top=2mm,bottom=2mm,before skip=4pt,after skip=4pt,
  fontupper=\small\RaggedRight
}

% Puente narrativo entre secciones (contrato editorial Conéctate).
% Da continuidad: "Acabas de X. Ahora vas a Y porque Z."
% Visualmente: banda izquierda en color del grado, texto en itálica pequeña.
\newtcolorbox{puentebox}{
  enhanced,colback=milcGradeSoft,colframe=milcGradeDark,
  borderline west={3pt}{0pt}{milcGradeDark},
  arc=0mm,boxrule=0pt,left=5mm,right=4mm,top=2.5mm,bottom=2.5mm,
  before skip=4pt,after skip=8pt,
  fontupper=\itshape\small\color{milcNegro}\RaggedRight
}
\newcommand{\puente}[1]{%
  \begin{puentebox}%
    \textcolor{milcGradeDark}{\textbf{\textrightarrow}}\hspace{2mm}#1%
  \end{puentebox}%
}

\newcommand{\linead}[1]{\noindent\color{milcLinea}\rule{#1}{0.5pt}\color{milcNegro}}
\newcommand{\respuesta}[1]{\par\vspace{2mm}\foreach \n in {1,...,#1}{\noindent\color{milcLinea}\rule{\linewidth}{0.5pt}\par\vspace{5mm}}\color{milcNegro}}
\newcommand{\checkbox}{\textcolor{milcGradeDark}{\fbox{\phantom{x}}}\hspace{5pt}}

\newcommand{\phasecard}[4]{
\begin{tcolorbox}[enhanced,colback=#3,colframe=#2,arc=3mm,boxrule=.5pt,height=3.05cm,valign=center,halign=center,left=1.5mm,right=1.5mm,top=2mm,bottom=2mm]
{\sffamily\bfseries\fontsize{22}{24}\selectfont\textcolor{#2}{#1}}\par
{\sffamily\bfseries\small #4}
\end{tcolorbox}
}

\begin{document}
\thispagestyle{empty}
\begin{tikzpicture}[remember picture,overlay]
  \fill[white] (current page.south west) rectangle (current page.north east);
  \path[fill=milcGradePrimary,rounded corners=16mm]
    ([xshift=.55cm,yshift=-.55cm]current page.north west) rectangle
    ([xshift=-.55cm,yshift=3.65cm]current page.south east);

  \node[anchor=north west,text=milcGradeText] at ([xshift=2.0cm,yshift=-1.85cm]current page.north west)
    {\sffamily\bfseries\fontsize{14}{16}\selectfont GRADO};
  \node[anchor=north west,text=milcGradeText,opacity=.82] at ([xshift=2.0cm,yshift=-2.75cm]current page.north west)
    {\sffamily\bfseries\fontsize{9}{11}\selectfont SERIE GUÍAS · TECNOLOGÍA E INFORMÁTICA};

  \begin{scope}[shift={([xshift=-3.05cm,yshift=-2.55cm]current page.north east)}]
    \fill[white,opacity=.80] (-.62,-.52) rectangle (.62,.52);
    \draw[milcGradeText,line width=1pt] (-.62,-.52) rectangle (.62,.52);
    \fill[milcGradeDark] (-.42,-.38) rectangle (-.22,.06);
    \fill[milcGradeAccent] (-.10,-.38) rectangle (.10,.24);
    \fill[milcGradePrimary!60!black] (.22,-.38) rectangle (.42,.38);
  \end{scope}

  \node[anchor=west,text=milcGradeText] at ([xshift=1.9cm,yshift=-12.85cm]current page.north west)
    {\sffamily\bfseries\fontsize{124}{124}\selectfont 7};
  \draw[milcGradeText,line width=1.7pt]
    ([xshift=4.52cm,yshift=-11.68cm]current page.north west) circle (.13cm);

  \node[anchor=west,text=milcGradeText,text width=8.7cm,align=left] at ([xshift=10.35cm,yshift=-11.6cm]current page.north west)
    {
     {\sffamily\bfseries\fontsize{19}{22}\selectfont Coautoría\\asincrónica ---\\3 personas\\al mismo\\tiempo}\\[4mm]
     {\sffamily\bfseries\fontsize{23}{26}\selectfont Guía 6-7-TIC}\\[2mm]
     {\sffamily\fontsize{13}{16}\selectfont Período 1 · Trabajo colaborativo en la nube}\\[5mm]
     {\sffamily\bfseries\fontsize{11}{13}\selectfont Institución Educativa Sor María Juliana}\\[1mm]
     {\sffamily\fontsize{10}{12}\selectfont Autor: Dr. Álvaro Cárdenas Orozco}
    };

  \path[fill=milcGradeSoft,rounded corners=7mm]
    ([xshift=1.35cm,yshift=.85cm]current page.south west) rectangle
    ([xshift=-1.35cm,yshift=4.25cm]current page.south east);
  \draw[milcGradeDark,line width=.65pt,rounded corners=7mm]
    ([xshift=1.35cm,yshift=.85cm]current page.south west) rectangle
    ([xshift=-1.35cm,yshift=4.25cm]current page.south east);
  \node[anchor=center,text=milcNegro,text width=17.0cm,align=center] at ([yshift=2.55cm]current page.south)
    {\sffamily\fontsize{10}{12}\selectfont Metodología MILC · Apertura ancestral · Pensamiento computacional · Triángulo Dussel-Estoicismo-Floridi\\
    \textcolor{milcGradeDark}{\bfseries Producto:} Un \textbf{documento Word colaborativo} entre 3 estudiantes (tú + 2 compañeros). El documento tiene: \textbf{(a) acuerdo de roles} arriba (quién escribe qué sección); \textbf{(b) 3 secciones distintas} con encabezados claros (una por estudiante); \textbf{(c) tu sección personal} con 1 párrafo de 4-5 líneas sobre un subtema del tema acordado. Más \textbf{en el cuaderno}: bitácora de lo que pasó (qué funcionó, qué se complicó, cómo lo resolvieron).};
\end{tikzpicture}
\mbox{}
\newpage

\section*{Apertura: Coautoría asincrónica --- 3 personas escribiendo al mismo tiempo}

\begin{softbox}{milcGradeDark}{milcGradeSoft}{Saber ancestral}
\textbf{En una minga grande del Valle del Cauca, los 30 trabajadores no estaban todos haciendo lo mismo.} En la finca cafetera de Cartago, cuando llegaba la minga, había \textbf{división espontánea de tareas}: \emph{10 personas en los cafetales recogiendo}, \emph{5 en los patios extendiendo el grano al sol}, \emph{3 en las cocinas preparando almuerzo}, \emph{2 cargando agua del nacedero}, \emph{1 cuidando que los niños no se metieran al cafetal}. Aunque eran 30 trabajando \textbf{al mismo tiempo}, cada quien sabía \emph{exactamente qué hacer y dónde}. \textbf{No había choques}. Esa organización no salía del aire: el dueño de la finca y los líderes de la minga la habían \emph{acordado al inicio del día}, mientras tomaban el café del desayuno. \textbf{Sin ese acuerdo inicial, los 30 se hubieran tropezado entre sí}. Hoy en Microsoft 365, la coautoría real --- 3 o más personas escribiendo en el mismo Word al tiempo --- funciona igual: \emph{primero acuerdo de roles, después trabajo paralelo, sin choques}.
\end{softbox}

\begin{softbox}{milcTurquesa}{milcGris}{Saber contemporáneo}
\textbf{Coautoría} es cuando varias personas escriben en el mismo archivo \textbf{a la vez}, viendo los cambios de los otros en tiempo real. Microsoft 365 (Word, Excel, PowerPoint en línea) permite hasta \textbf{cientos de coautores simultáneos}. Cuando varios escriben al mismo tiempo, ves los \emph{cursores de los demás} en colores distintos, con sus nombres flotando. \textbf{Asincrónica} significa que NO requiere que todos estén conectados al mismo tiempo: cada uno entra cuando puede, escribe su parte, sale; el documento se actualiza en la nube. \textbf{Sincrónica} es cuando todos están conectados a la vez (más rápido, pero más caótico si no hay coordinación). Para 7° grado lo ideal es \textbf{asincrónica con acuerdo previo}: definen secciones, cada uno trabaja en la suya, y se reúnen al final para revisar.
\end{softbox}

\begin{softbox}{milcMostaza}{milcGris}{Pregunta-puente}
\textit{Si tú y 2 compañeros tuvieran que escribir un trabajo de Sociales de 6 párrafos en 1 hora trabajando \textbf{cada uno desde su casa}, ¿\textbf{cómo se organizarían}? ¿Cómo evitan escribir cada uno la misma cosa o, peor, borrarse entre sí?}
\end{softbox}

\begin{softbox}{milcVerde}{milcGris}{Saber y saber hacer}
Al terminar la clase: (1) podrás \textbf{identificar} las diferencias entre coautoría sincrónica y asincrónica; (2) sabrás \textbf{aplicar} los 4 pasos del acuerdo de roles; (3) podrás \textbf{evaluar} si una coautoría está bien organizada; (4) habrás \textbf{creado} un documento colaborativo con 2 compañeros.
\end{softbox}

\small
\begin{tabularx}{\textwidth}{>{\bfseries}p{3.0cm}Y}
\rowcolor{milcGradePrimary!12}Autor & Dr. Álvaro Cárdenas Orozco\\
DBA & Utiliza herramientas digitales para trabajar de manera colaborativa con otros (MEN, T\&I 7°)\\
Referentes & Saber ancestral de la minga · Microsoft 365 y OneDrive · Coautoría asincrónica\\
Producto final & Un \textbf{documento Word colaborativo} entre 3 estudiantes (tú + 2 compañeros). El documento tiene: \textbf{(a) acuerdo de roles} arriba (quién escribe qué sección); \textbf{(b) 3 secciones distintas} con encabezados claros (una por estudiante); \textbf{(c) tu sección personal} con 1 párrafo de 4-5 líneas sobre un subtema del tema acordado. Más \textbf{en el cuaderno}: bitácora de lo que pasó (qué funcionó, qué se complicó, cómo lo resolvieron).\\
\end{tabularx}
\normalsize

\puente{Hoy haces 4 cosas: descubres los 4 pasos del acuerdo, formas equipo de 3, escriben el documento, registran la experiencia.}

\begin{titlebox}{milcGradeDark}
Ruta MILC: así vamos a aprender
\end{titlebox}
\begin{tabularx}{\textwidth}{>{\centering\arraybackslash}X>{\centering\arraybackslash}X>{\centering\arraybackslash}X>{\centering\arraybackslash}X}
\phasecard{1}{milcVerde}{milcGris}{Escuta\\[-1mm]\small Observo, escucho y pregunto.} &
\phasecard{2}{milcTurquesa}{milcGris}{Sistematización\\[-1mm]\small Pensamiento computacional.} &
\phasecard{3}{milcMagenta}{milcGris}{Praxis\\[-1mm]\small Construyo y aplico.} &
\phasecard{4}{milcVino}{milcGris}{Evaluación\\[-1mm]\small Reflexión liberadora.}
\end{tabularx}


\puente{Empezamos comparando coautoría buena vs coautoría caótica.}

\begin{titlebox}{milcVerde}
Fase 1 · Escuta
\end{titlebox}

\begin{softbox}{milcVerde}{milcGris}{Escena para escuchar}
\textbf{Actividad 1 · ANALIZA --- 2 historias de coautoría} (10 min · individual). Lee estas 2 historias y en el cuaderno responde: \emph{¿qué hizo distinta a la historia B de la A?}. \textbf{Historia A (caos):} \emph{``María, Pedro y Lucía tienen que entregar un Word de 6 párrafos sobre los ríos de Colombia. Se conectan al documento al mismo tiempo, sin haber acordado nada. Los 3 empiezan a escribir desde la primera línea. Se borran entre sí. María se enoja porque Pedro borró lo que ella había escrito. Lucía siente que está sobrando. A la hora se rinden y deciden cada uno hacer su versión por separado. Pierden 2 horas.''}. \textbf{Historia B (éxito):} \emph{``María, Pedro y Lucía abren WhatsApp. María propone: yo hago río Magdalena, Pedro río Cauca, Lucía río Atrato. Se acuerdan los 6 párrafos: 2 por río. Cada uno entra al Word, busca su sección y escribe sin pisarse. En 45 minutos terminan. Se reúnen 10 minutos a revisar y entregar.''}. \textbf{Responde:} ¿qué hizo posible la historia B?
\end{softbox}

\checkbox Leí las 2 historias con calma.\\
\checkbox Identifiqué la diferencia clave.\\
\checkbox Pensé en mi propia experiencia con trabajos en grupo.

\begin{infoband}{milcVerde}
\textbf{Lo que va al cuaderno (Actividad 1 · ANALIZA).} Título: «Actividad 1 --- 2 historias de coautoría». \textbf{Formato:} respuesta corta (3-4 líneas) sobre qué hizo distinta la B. \textbf{Extensión:} 4 líneas. \textbf{Sabes que terminaste cuando} mencionas el acuerdo previo como clave.
\end{infoband}

\puente{Después aprendes los 4 pasos del acuerdo de roles y las 5 técnicas para no chocar.}

\begin{titlebox}{milcTurquesa}
Fase 2 · Sistematización con pensamiento computacional
\end{titlebox}

\begin{softbox}{milcTurquesa}{milcGris}{Cuatro pilares aplicados al tema}
La diferencia entre la historia A (caos) y la historia B (éxito) es solo \textbf{1 cosa}: \emph{acuerdo de roles antes de empezar}. Ahora aprendes los 4 pasos exactos del acuerdo y las 5 técnicas para no chocar.
\end{softbox}

\small
\begin{tabularx}{\textwidth}{>{\bfseries\RaggedRight\arraybackslash}p{3.6cm}Y}
\rowcolor{milcTurquesa!90}\color{white}Pilar & \color{white}\bfseries Aplicación al tema\\
1. Descomposición & \textbf{Los 4 pasos del acuerdo de roles (5 minutos al inicio).} \textbf{(1) Lectura común del tema:} todos los 3 leen el enunciado del trabajo. \emph{``Hacer 6 párrafos sobre los ríos de Colombia''}. Sin lectura común, cada uno entiende algo distinto. \textbf{(2) Lluvia de subtemas:} entre los 3 listan 3-5 subtemas posibles. \emph{``Río Magdalena, río Cauca, río Atrato, río Orinoco''}. \textbf{(3) Asignación de roles:} cada uno toma un subtema. \emph{``Yo Magdalena, tú Cauca, ella Atrato''}. Si alguien quiere un tema específico, hablan y deciden. \textbf{(4) Tiempos y entregas:} acuerdan cuándo termina cada uno. \emph{``Yo termino mi parte el miércoles 6pm, tú jueves 8pm, ella viernes 10am. Revisamos juntos viernes 11am''}.\\
2. Reconocimiento de patrones & \textbf{Las 5 técnicas para no chocar durante la coautoría.} (1) \textbf{Secciones bien separadas:} usar encabezados grandes con el nombre del autor (\emph{``Río Magdalena --- María''}). (2) \textbf{Mirar antes de escribir:} cuando entres al documento, fíjate si hay otro cursor en tu zona. Si lo hay, espera un poco o escríbele un comentario. (3) \textbf{No tocar las secciones de los demás}: aunque veas un error o se te ocurra algo mejor. Esa sección es de ellos. (4) \textbf{Comunicarse en paralelo:} mantener un chat de WhatsApp o Teams abierto para coordinar dudas en tiempo real. (5) \textbf{Guardado de seguridad:} cuando termines tu parte, copia tu sección en un archivo aparte en tu propio OneDrive. Así, si algo se daña en el compartido, tu trabajo está respaldado.\\
3. Abstracción & \textbf{Las 3 herramientas visuales de Word para coautoría.} \textbf{(a) Cursores con nombres:} cuando otros escriben, ves dónde están con su cursor de color y su nombre flotando. Útil para saber quién está activo y dónde. \textbf{(b) Comentarios al margen:} si quieres sugerir algo en la sección de otro sin escribir directamente, agregas un comentario (clic derecho → Nuevo comentario). El otro lo ve y decide. \textbf{(c) Pestaña ``Personas'':} arriba a la derecha del documento, muestra quiénes están conectados al mismo tiempo (sus íconos). Útil para saber con quién estás coautorando ahora.\\
4. Algoritmo conceptual & \textbf{Errores típicos de equipos novatos en coautoría.} (1) \emph{Empezar sin acuerdo} --- la historia A. (2) \emph{Un solo encargado que hace todo y los demás miran} --- pierde el punto de la coautoría. (3) \emph{Todos escriben en el primer párrafo} --- se borran entre sí. (4) \emph{Nadie revisa al final} --- el documento queda con estilos distintos, fuentes diferentes, contradicciones. (5) \emph{Los 3 son responsables pero al final 1 hace el trabajo de los 2 que no aportaron} --- comunicación clara desde el inicio + responsabilidad lo evita.\\
\end{tabularx}
\normalsize

\begin{drawbox}{milcTurquesa}{4 pasos + 5 técnicas + 3 herramientas}
\textbf{4 pasos del acuerdo:} lectura común · lluvia de subtemas · asignación · tiempos. \\[1mm]
\textbf{5 técnicas para no chocar:} secciones separadas · mirar antes · respetar sección ajena · chat paralelo · backup propio. \\[1mm]
\textbf{3 herramientas Word:} cursores con nombres · comentarios al margen · pestaña Personas.
\end{drawbox}

\begin{softbox}{milcMostaza}{milcGris}{Errores comunes y su corrección}
\textbf{Mitos sobre la coautoría que te encontrarás.} (1) \emph{``Coautoría es muy difícil, mejor cada uno hace su parte aparte''} --- Falso. Con acuerdo claro es más rápido. (2) \emph{``Si soy el mejor, hago todo y firmamos los 3''} --- Mala idea: tus compañeros no aprenden y tú haces trabajo extra. (3) \emph{``En coautoría los cambios se pierden si todos escriben a la vez''} --- Falso. Word maneja muy bien los cambios simultáneos siempre que estén en secciones distintas. (4) \emph{``Coautoría requiere que todos estén conectados al mismo tiempo''} --- Falso. La asincrónica permite que cada uno trabaje cuando pueda.
\end{softbox}

\begin{infoband}{milcTurquesa}
\textbf{Lo que va al cuaderno (Actividad 2 · EXPLICA).} Título: «Actividad 2 --- 4 pasos + 5 técnicas». \textbf{Formato:} 2 bloques. Bloque A: los 4 pasos numerados con 1 línea cada uno. Bloque B: las 5 técnicas numeradas con 1 línea cada una. \textbf{Extensión:} 9 líneas. \textbf{Sabes que terminaste cuando} podrías guiar a 2 compañeros nuevos en coautoría.
\end{infoband}

\puente{Con todo claro, formas equipo, acuerdan roles, escriben en paralelo.}

\begin{titlebox}{milcMagenta}
Fase 3 · Praxis con pensamiento computacional
\end{titlebox}

\begin{softbox}{milcMagenta}{milcGris}{Construyendo el producto con los cuatro pilares}
\textbf{Actividad 3 · CREA --- Documento colaborativo con 2 compañeros} (35 min · equipo de 3). Esta actividad es \textbf{en equipo}: tú + 2 compañeros van a producir un Word colaborativo en tiempo real (en clase si tienen computadores; o asincrónico durante la semana si no). Tema sugerido: \emph{``3 ríos importantes de Colombia''} o \emph{``3 platos tradicionales del Valle del Cauca''} o lo que acuerden.
\end{softbox}

\small
\begin{tabularx}{\textwidth}{>{\bfseries\RaggedRight\arraybackslash}p{3.6cm}Y}
\rowcolor{milcMagenta!90}\color{white}Pilar & \color{white}\bfseries Aplicación al producto\\
Descomposición & \textbf{Paso 1 --- Formen equipo y abran chat (3 min).} Únanse con 2 compañeros. Abran un chat en WhatsApp o Teams. Definan nombre del equipo. Saluden y propongan tema (de los sugeridos o uno propio).\\
Patrones & \textbf{Paso 2 --- Acuerdo de roles (5 min).} Sigan los 4 pasos: \emph{(a) lectura común del enunciado}, \emph{(b) lluvia de subtemas} (escriban 3 subtemas en el chat), \emph{(c) asignación} (cada uno toma uno), \emph{(d) tiempos} (acuerden cuándo termina cada uno). Anótenlo en el chat para que nadie olvide.\\
Abstracción & \textbf{Paso 3 --- Uno de los 3 crea el documento y comparte como editor (3 min).} El que tenga más signal/wifi rápido abre Word Online, crea \emph{Documento en blanco}, lo nombra \texttt{2026-trabajo-colaborativo-[tema]}. En la parte superior escribe el \textbf{acuerdo de roles} (los 3 subtemas + los 3 autores + los 3 tiempos acordados). Después comparte con los otros 2 con permiso \textbf{Puede editar}.\\
Algoritmo de armado & \textbf{Paso 4 --- Escriben cada uno su sección (15-20 min).} Cada uno entra al documento desde su computador (o celular si toca). Busca su sección por su encabezado. \textbf{Importante:} crea primero \textbf{los 3 encabezados grandes} (uno por subtema con nombre del autor). Cada uno escribe \textbf{1 párrafo de 4-5 líneas} sobre su subtema. Aplica formato (negrita en ideas clave, justificado). \emph{Si entras al documento y ves que tu compañero ya está en otra sección, tranquilo: tienen secciones distintas, no se chocan}.\\
\end{tabularx}
\normalsize

\begin{drawbox}{milcMagenta}{Antes de cerrar la sesión 6}
\checkbox Formé equipo con 2 compañeros y acordamos roles. \\[1mm]
\checkbox El documento tiene 3 secciones con encabezados claros. \\[1mm]
\checkbox Cada estudiante escribió SU sección sin tocar las otras. \\[1mm]
\checkbox El documento se ve uniforme (mismo formato general). \\[1mm]
\checkbox Mi bitácora individual tiene las 6 preguntas respondidas. \\[1mm]
\checkbox La bitácora está firmada con nombre y fecha.
\end{drawbox}

\begin{softbox}{milcMagenta}{milcGris}{Plantilla de trabajo}
\textbf{Estructura del trabajo de equipo.} \textit{Paso 1:} formar equipo + chat + tema. \textit{Paso 2:} acuerdo de 4 pasos. \textit{Paso 3:} crear y compartir documento. \textit{Paso 4:} escribir cada uno su sección. \textit{Paso 5:} revisión conjunta. \textit{Paso 6:} bitácora individual.
\end{softbox}

\begin{infoband}{milcMagenta}
\textbf{Lo que va al cuaderno (Actividad 3 · CREA).} Título: «Actividad 3 --- Mi primer documento colaborativo». \textbf{Formato:} datos del equipo + acuerdo de roles + bitácora individual de 6 puntos. \textbf{Extensión:} 1 página. \textbf{Sabes que terminaste cuando} el documento colaborativo está completo Y tu bitácora está firmada.
\end{infoband}

\puente{El producto es el documento colaborativo + bitácora.}

\begin{titlebox}{milcMostaza}
Producto de la sesión
\end{titlebox}
\textbf{Documento Word colaborativo de 3 secciones + bitácora individual · firmada}.
\begin{softbox}{milcMostaza}{milcGris}{Criterios de calidad}
(1) El documento tiene acuerdo de roles visible arriba. (2) Hay 3 secciones distintas con encabezados claros. (3) Cada autor escribió un párrafo de 4-5 líneas en su sección. (4) El documento mantiene formato uniforme (no es 'collage' de estilos distintos). (5) La bitácora individual tiene 6 puntos respondidos con honestidad.
\end{softbox}

\puente{Antes de salir, miran el documento juntos: ¿se ve como uno solo o como 3 pegados?}

\begin{titlebox}{milcVino}
Fase 4 · Evaluación liberadora — cinco dimensiones
\end{titlebox}

\begin{softbox}{milcVino}{milcGris}{Sentido de la evaluación}
Evaluar no es solo revisar una nota. Es reconocer qué aprendí, qué cambió en mí, qué emociones manejé, cómo me relaciono con la ciudad y la comunidad, y qué le dejaré a quien viene detrás.
\end{softbox}

\textbf{Mi reflexión liberadora en cinco dimensiones.} Completa cada frase iniciada:

\small
\begin{tabularx}{\textwidth}{>{\bfseries\RaggedRight\arraybackslash}p{3.4cm}Y>{\RaggedRight\arraybackslash}p{4.6cm}}
\rowcolor{milcVino}\color{white}Dimensión & \color{white}\bfseries Mi respuesta (frame starter) & \color{white}\bfseries Algo que descubrí\\
1. Personal & Lo que más cambió en mí fue \linead{6.5cm} & \linead{4.2cm}\\
2. Emocional & La emoción más fuerte fue \linead{2.5cm} y la regulé \linead{3cm} & \linead{4.2cm}\\
3. Ciudadana & En democracia esto importa porque \linead{6cm} & \linead{4.2cm}\\
4. Local (Cartago) & En mi barrio sería útil para \linead{6.5cm} & \linead{4.2cm}\\
5. Intergeneracional & A mi abuela/abuelo le diría \linead{6.5cm} & \linead{4.2cm}\\
\end{tabularx}
\normalsize

\begin{infoband}{milcVino}
\textbf{Para la próxima sustentación.} Vuelve a esta tabla en seis meses. Verás que las respuestas cambian — no porque lo aprendido cambie, sino porque tú cambias. La evaluación liberadora es una conversación contigo a través del tiempo.
\end{infoband}


\puente{Cerramos con 3 voces sobre el oficio de trabajar juntos sin chocar.}

\begin{titlebox}{milcGradeDark}
Triángulo de pensamiento · cierre filosófico
\end{titlebox}

\begin{softbox}{milcVino}{milcGris}{Dussel · lente del nosotros}
\textit{````No hay obra grande sin reparto previo. Quien no acuerda al inicio, choca al final.''''}

La diferencia entre la historia A (caos) y la historia B (éxito) no estuvo en el talento de los compañeros, ni en la herramienta, ni en la suerte. Estuvo en \textbf{5 minutos de acuerdo al inicio}. La minga lo sabía hace 100 años. Dussel lo dice en filosofía. Tu equipo lo vive en práctica. \emph{Acordar al inicio ahorra horas al final}. Es una habilidad que vale toda la vida.

\textbf{¿Cuántos trabajos de grupo se me han complicado por falta de acuerdo inicial?}
\end{softbox}
\linead{\linewidth}\\[1mm]
\linead{\linewidth}

\begin{softbox}{milcMostaza}{milcGris}{Estoicismo · Marco Aurelio · cuidado interior}
\textit{````Confía en tu equipo lo suficiente para no querer hacer su parte. Esa confianza es la base del trabajo colectivo.''''}

Hay una tentación grande en coautoría: \emph{``yo lo hago todo porque sé que sale mejor''}. Mala idea. Cuando le quitas la parte al compañero, le quitas también el aprendizaje. Y te cargas tú con trabajo de 3. \textbf{Confiar en el equipo es soltar el control sin abandonar la responsabilidad}. Esa habilidad es de adulto. Empezar a desarrollarla en 7° grado es ventaja real.

\textbf{¿Confío en mi equipo o termino haciendo todo yo por desconfianza?}
\end{softbox}
\linead{\linewidth}\\[1mm]
\linead{\linewidth}

\begin{softbox}{milcTurquesa}{milcGris}{Floridi · lente de la infoesfera}
\textit{````La coautoría digital es la primera escuela de ciudadanía colectiva. Ahí aprendes a hacer cosas grandes con otros.''''}

Los grandes proyectos del mundo --- Wikipedia, GitHub, las series de Netflix, los reportajes de El Tiempo --- los hacen \textbf{equipos de cientos de personas coautorando}. Aprender a coautorar bien en 7° grado te conecta con esa forma de producir. \emph{Es la habilidad central de la economía del conocimiento del siglo XXI}. No es tema escolar: es preparación profesional.

\textbf{¿Estoy desarrollando habilidad para hacer cosas grandes con otros, o me quedo en el modo individual?}
\end{softbox}
\linead{\linewidth}\\[1mm]
\linead{\linewidth}

\begin{titlebox}{milcVino}
Compromiso final
\end{titlebox}

\small
\begin{tabularx}{\textwidth}{>{\RaggedRight\arraybackslash}p{3.0cm}YYY}
\rowcolor{milcVino}\color{white}\bfseries Criterio & \color{white}\bfseries Lo logré & \color{white}\bfseries En proceso & \color{white}\bfseries Necesito apoyo\\
Comprensión del tema & Explico ideas centrales con mis palabras. & Reconozco ideas pero falta precisión. & Necesito repasar conceptos base.\\
Pensamiento computacional & Apliqué los 4 pilares al diseñar mi producto. & Apliqué algunos pilares. & Necesito releer la fase 2.\\
Aplicación y producto & Mi producto cumple los criterios y se puede revisar. & Producto funcional con ajustes pendientes. & Aún en construcción.\\
Reflexión 5 dimensiones & Respondí cada dimensión con honestidad. & Respondí algunas con detalle. & Mis respuestas son muy generales.\\
Triángulo de pensamiento & Conecté las 3 voces con mi trabajo real. & Conecté al menos una. & No relaciono las voces con mi proceso.\\
\end{tabularx}
\normalsize

\textbf{Hoy entendí que:}\\[1mm]
\linead{\linewidth}\\[1mm]
\linead{\linewidth}

\textbf{Mi compromiso para la próxima sesión es:}\\[1mm]
\linead{\linewidth}\\[1mm]
\linead{\linewidth}

\begin{softbox}{milcMostaza}{milcGris}{Cita de despedida · Marco Aurelio}
\textit{``El obstáculo en el camino se vuelve el camino. Lo que estorba al camino se vuelve camino.''}\\[1mm]
Cada error de hoy, bien observado, se convierte en pieza del camino que pisarás mañana.
\end{softbox}

\small
\begin{tabularx}{\textwidth}{>{\bfseries}p{3.6cm}Y}
\rowcolor{milcGradeDark!12}Nombre completo: & \linead{10cm}\\
Fecha: & \linead{4cm} \quad Curso/grupo: \linead{4cm}\\
Mi firma: & \linead{9cm}\\
\end{tabularx}
\normalsize

\begin{infoband}{milcGradeDark}
\textbf{Para la próxima sesión.} Esta semana, revisa el documento colaborativo con tus 2 compañeros y haz ajustes finales si se necesitan. La próxima clase aprendes a \textbf{comentar y sugerir} sin escribir directamente --- el oficio de revisar el trabajo ajeno con respeto. Es lo que falta para que tu equipo se vuelva profesional.
\end{infoband}

\end{document}
